\section{Conditional density and field comparison at positive price}
\label{sec:comparison}

We now obtain a local Thomas--Fermi relation for a conditional density.
The conditioning retains information about the electron configuration while
permitting a controlled change of state on any event of appreciable probability.
Two different conditional laws enter the argument: observations define the
density to be estimated, whereas a fresh spatial cut defines the core used in
the energy comparison.  We keep these laws separate and match their potentials
only after taking expectation.

Throughout this section, $\Psi$ is a normalized ground state in any sector $n$
minimizing $P(\lambda)$, with $\lambda=s^{-4}$ and
\cref{eq:scaling-regime} in force.  In particular,
$n\leq 3Z$ and $D_0=E_n-E\leq Cs^{-7}$ by
\cref{lem:priced-rough}.  All statements about sufficiently small scales are
uniform in the choice of this minimizing sector.

\subsection{Nested observations and their energy cost}

Fix $0<\eps<1$, to be chosen in \cref{sec:barriers}, and define
\begin{equation}
 r_0=\eps Z^{-1/3},\qquad r_j=2^j r_0,\qquad
 J_* =\max\{j\geq0:r_j\leq s\},\qquad
 \ell_j=r_j^{1.01}\quad(0\leq j<J_*).
 \label{eq:observation-scales}
\end{equation}
Eventually $2r_0\leq s$.  At each scale $j<J_*$, observe the unordered array
of points $x_i+\ell_jU_{ji}$.  All coordinates of all $U_{ji}$ are independent,
also independently of the raw configuration, with density
\[
 \zeta(u)=c_\zeta\exp\!\left(-\frac1{1-u^2}\right)
                  \ind_{\{|u|<1\}}.
\]
Labels are forgotten separately in each array.  Let $\mathcal F_j$ be the
information in arrays $j,\ldots,J_*-1$, and let $\mathcal F_{J_*}$ be trivial.
These sigma-fields decrease with $j$.

Use the constants $L=10^5$, $A=40$ and $a_y=|y|/L$ from
\cref{lem:local-field-input}.  Fix a real radial function
$g\in C_c^\infty(B(0,1))$ with $\int g^2=1$, set
$g_b(x)=b^{-3/2}g(x/b)$, and choose
\[
 w=10^{-5},\qquad 0<c_1<(10L)^{-1}.
\]
Writing $d_x=|x|$, define the master packets and conditional fields by
\begin{equation}
 \begin{aligned}
 D_x^{\mathrm b}&=\max(d_x,r_0),&
 t_x&=c_1D_x^{\mathrm b}\min(D_x^{\mathrm b},s)^w,\\
 \mathcal K_x(y)&=g_{t_x}(y-x)^2,&
 \mu_j(y)&=\E\!\left[\sum_i\mathcal K_{x_i}(y)
                                  \,\middle|\,\mathcal F_j\right],\\
 H'_j(y)&=V(y)-\lambda-(K*\mu_j)(y).
 \end{aligned}
 \label{eq:master-kernel}
\end{equation}
The superscript $\mathrm b$ indicates the lower clamp, not a power.
The width function is globally Lipschitz, with
$\operatorname{Lip}(t)\leq c_1(1+w)s^w$, and
$t_x\geq c_1r_0^{1+w}$.  Thus, for each fixed system, $\mu_j$ is smooth in
its evaluation variable, has mass $n$, and has deterministic bounds on all
spatial derivatives.  Its Coulomb potential and first derivatives are also
continuous and deterministically bounded.  These bounds at a fixed system
will justify conditional integration; they are not claimed to be uniform in
$Z$.

We use versions obtained by integrating against the conditional law of the
raw configuration.  Such versions can, for example, be specified by the
finite-array likelihoods in \cite[\FObservations]{Foundation}.  They give,
simultaneously in space,
\begin{equation}
 \E[\mu_j\mid\mathcal F_{j+1}]=\mu_{j+1},\qquad
 \E[H'_j\mid\mathcal F_{j+1}]=H'_{j+1},\qquad
 \Delta H'_j=-4\pi\nu+4\pi\mu_j.
 \label{eq:observation-tower}
\end{equation}
Continuity extends identities first obtained on a countable dense set.
In particular, $H'_j$ is subharmonic off the nucleus.

\begin{lemma}[Cost of an observation event]
\label{lem:event-tilt}
Let $j<J_*$ and let $A'\in\mathcal F_j$ have probability $p>0$.
There is a normalized antisymmetric state $\Psi_{A'}$ in the same sector
whose raw law is the raw marginal conditioned on $A'$ and whose form energy
satisfies
\begin{equation}
 \langle\Psi_{A'},H_{Z,n}\Psi_{A'}\rangle
 \leq E_n+C r_j^{-2.02}\log^5(e/p).
 \label{eq:event-tilt-cost}
\end{equation}
The constant is independent of the particle number, the number of observed
arrays, and $D_0$.
\end{lemma}

\begin{proof}
If $p_{A'}(x)$ is the probability of the observed event conditional on the
raw positions, take
$\Psi_{A'}=\sqrt{p_{A'}(x)/p}\,\Psi$.
The multiplier is symmetric and preserves the form domain.  The first
inequality of \cite[\FEventTilt]{Foundation}, with
$\ell_j=r_j^{1.01}$, is exactly \cref{eq:event-tilt-cost}.
Its proof uses the ground-state multiplier identity in the sector $n$;
it measures the cost above $E_n$, independently of the offset of $E_n$
from $E$.
\end{proof}

Accordingly, the unshifted offset of the tilted state is at most
$D_0+C r_j^{-2.02}\log^5(e/p)$, while its shifted energy is at most
\begin{equation}
 P(\lambda)+C r_j^{-2.02}\log^5(e/p).
 \label{eq:shifted-tilt-cost}
\end{equation}
This distinction is what permits the local comparison below for
$D_0$ up to order $s^{-7}$.

\subsection{Statement and preliminary bounds}

\begin{proposition}[Simultaneous inverse comparison]
\label{prop:inverse-comparison}
There is a universal constant $C_{\mathrm{cap}}$ with the following property.
Fix
\[
 M\geq2,\qquad 0<h_l<h_h<\infty,\qquad 0<\xi<h_l/16.
\]
For all sufficiently far systems along \cref{eq:scaling-regime}, at every
$j<J_*$ there is an event $\mathcal G_j\in\mathcal F_j$ such that, with
$r=r_j$,
\[
 \Prob(\mathcal G_j^c)\leq Cr^{25}.
\]
On $\mathcal G_j$, the following assertions hold simultaneously for
$r\leq d_y\leq4Mr$ and $h_l\leq h\leq h_h$:
\begin{equation}
 \begin{split}
 d_y^4H'_j(y)&\leq C_{\mathrm{cap}},\\
 d_y^6\mu_j(y)>kh^{3/2}
       &\quad\Longrightarrow\quad d_y^4H'_j(y)\geq h-\xi,\\
 d_y^6\mu_j(y)<kh^{3/2}
       &\quad\Longrightarrow\quad d_y^4H'_j(y)\leq h+\xi.
 \end{split}
 \label{eq:inverse-comparison}
\end{equation}
The probability constant and the threshold for sufficiently far systems
may depend on the displayed fixed parameters and on the fixed packet
conventions; $C_{\mathrm{cap}}$ is independent of the band and accuracy
parameters.
\end{proposition}

The two implications compare the field with each positive density height
from opposite sides.  They will yield an approximate Thomas--Fermi relation
in \cref{sec:limit}, while permitting the field itself to be arbitrarily
negative where the density is small.

We prove the proposition in the rest of this section.  The proof uses the
arbitrary-offset screening estimates from \cref{sec:pricing} and the energy
cost in \cref{lem:event-tilt}; it does not apply the fixed-offset inverse
comparison of \cite[\FFixedOffset]{Foundation} with a growing $D_0$.

Fix $j<J_*$ and put $r=r_j$.  On the enlarged band
$r/2\leq d_y\leq8Mr$, write
\[
 a=a_y,\qquad t=t_y,\qquad \tau=t/a.
\]
Uniformly there, for fixed $M$,
\begin{equation}
 a\asymp_M r,\qquad c_Ma^w\leq\tau\leq Cs^w.
 \label{eq:packet-scale-ratios}
\end{equation}
If $\mathcal K_x(y)\ne0$, the Lipschitz bound for the widths gives
$|x-y|\leq2t$ and $t_x/t=1+O(s^w)$.  Changing variables at scale $t$
therefore gives
\begin{equation}
 \int\mathcal K_z(y)\,dz=1+O(s^w).
 \label{eq:packet-approximate-mass}
\end{equation}
On the local support, the kernel is bounded by $Ct^{-3}$ and is Lipschitz,
in either its center or its evaluation variable, with constant $Ct^{-4}$.

\paragraph{A count event.}
Except on an event of probability at most $r^{40}$, the $j$th observed
array has at most $C_Mr^{-3}$ entries in the radial band
$[r/4,12Mr]$.  To prove this, condition on the event that a sufficiently
large such threshold is exceeded, if its probability is greater than
$r^{40}$.  Compact noise support forces every compatible raw
configuration to have at least that many points in $[r/8,13Mr]$.
By \cref{lem:event-tilt,eq:annular-count}, its conditional expected count
is at most $C_Mr^{-3}$: indeed,
\[
 \sqrt{D_0r}\leq Cr^{-3}(r/s)^{7/2}\leq Cr^{-3},\qquad
 r^{-0.51}\log^{5/2}(e/r^{40})=o(r^{-3}).
\]
This is a contradiction for a sufficiently large fixed threshold.
Denote the retained event by $\mathcal C_j$.  On it every compatible raw
configuration has at most $C_Mr^{-3}$ centers relevant to the enlarged
band.  Integration against the observation posterior consequently gives
\begin{equation}
 0\leq\mu_j(y)\leq C_Mr^{-6-3w},\qquad
 |\nabla\mu_j(y)|\leq C_Mr^{-7-4w}.
 \label{eq:conditional-density-bounds}
\end{equation}

\paragraph{A universal field cap.}
At a fixed point $y$ of the enlarged band, condition on
$\{d_y^4H'_j(y)>b_0\}$ if it has probability $p>0$.
Packets centered outside the exclusion ball in $J_y$ do not contain $y$,
by the width bound and the fixed choice of $c_1$.  Their potentials at
$y$ equal the corresponding point potentials by radial averaging.
Dropping the other nonnegative packet potentials and using the observation
tower therefore bounds the event-mean of $H'_j(y)$ above by the tilted
mean of $J_y-\lambda$.  The trivial-data case of
\cref{lem:local-field-input} gives
\begin{equation}
 b_0\leq C+C(d_y/s)^{7/2}-(d_y/s)^4
              +C_Mr^{2.49}\log^{5/2}(e/p).
 \label{eq:field-cap-tail}
\end{equation}
Here $a<1$ eventually.  The polynomial-power difference on the right
has a universal upper bound.  Choose a universal $C'$ exceeding the
entire deterministic bound by one.  Solving
\cref{eq:field-cap-tail} for $p$ and integrating its stretched-exponential
tail yields, for every fixed $P>0$,
\begin{equation}
 \E\bigl(d_y^4H'_j(y)-C'\bigr)_+
                    \leq C_{P,M}r^P.
 \label{eq:cap-excess}
\end{equation}
Subharmonicity turns this into a simultaneous bound.  For
$3r/4\leq d_y\leq6Mr$, take a ball of radius $d_y/10$; it lies in the
enlarged band.  Its mean bounds $H'_j(y)$ above.  The contribution of
$C'd_z^{-4}$, after multiplication by $d_y^4$, is bounded by a universal
constant, and the remaining contribution is at most
\[
 Cr^{-3}\int_{\{r/2\leq d_z\leq8Mr\}}
                         (d_z^4H'_j(z)-C')_+\,dz.
\]
By \cref{eq:cap-excess} and Markov's inequality, the latter is at most one
outside an event of probability $C_Mr^{40}$.  We have proved the
simultaneous cap
\begin{equation}
 H'_j(y)\leq C_{\mathrm{cap}}d_y^{-4}
       \quad(3r/4\leq d_y\leq6Mr)
 \label{eq:enlarged-field-cap}
\end{equation}
with a universal $C_{\mathrm{cap}}$.

\subsection{A fine patch comparison with the price included}

Fix a point $y$ in the original band and a height
$h_l/2\leq h\leq2h_h$.  Suppose that one of the following events,
intersected with $\mathcal C_j$, has probability $p\geq r^{42}$:
\begin{equation}
 \begin{array}{ll}
 d_y^6\mu_j(y)>kh^{3/2},&d_y^4H'_j(y)<h-\xi/4,\\[2pt]
 d_y^6\mu_j(y)<kh^{3/2},&d_y^4H'_j(y)>h+\xi/4.
 \end{array}
 \label{eq:fixed-point-failures}
\end{equation}
Write $Q$ for the original joint law conditioned on the event.
When spatial cuts are introduced, enlarge this law by fresh labels with
their ordinary conditional probabilities given the raw positions.
Its raw marginal is the state from \cref{lem:event-tilt}.  Set
\[
 \delta'=Cr^{-2.02}\log^5(e/r^{42}),\qquad
 \delta=D_0+\delta'.
\]
Then the raw state has unshifted offset at most $\delta$, shifted offset
at most $\delta'$, and
\begin{equation}
 m_y\leq C_Ma^{-3},\qquad
 \delta'=o(a^{-6.98}).
 \label{eq:patch-offsets}
\end{equation}
The constants below may depend on the fixed band parameters.

Put $\beta=a^{0.2}$ and $b=\beta a$, with $s$ small enough that
$\beta<1/100$.  For each fixed band parameter $M$, the scales satisfy
\[
 b\ll\ell_j\ll t\ll a.
\]
Here each ratio of consecutive scales tends to zero uniformly in $j$,
also when $r/s\to0$, by \cref{eq:packet-scale-ratios}.
The fine width $b$ makes the semiclassical errors small; the broader master
width $t$ absorbs both this smearing and the observation displacement
$\ell_j$, while $t/a\to0$ permits a local field comparison.
The patch construction below has three outputs.  Its energy comparison
controls the Coulomb discrepancy and a negative-field penalty.  The former,
together with noisy-array matching, recovers the conditional density; the
latter permits the high-density implication to pass to expectation.
Finally, the two conditional-expectation identities recover the mean
potential.  We establish these transfers in that order.

Choose a square cut, using the sine and cosine of a
smooth angle, that is full out on $B(y,t_0)$, supported out in
$B(y,t_0+b)$, and has gradients bounded by $C/b$, where
$5a\leq t_0\leq6a$ is selected below.  Retain the out particles satisfying
$|x_i-y|<t_0-7b$, and call the other out particles deleted.
Let $n_o$ denote the total number of out particles and
$n_{\mathrm{del}}$ the number of deleted particles; thus the retained
count is $n_o-n_{\mathrm{del}}$.  Then $t_0$ may be chosen so that
\begin{equation}
 \E_Q n_o^2\leq C_Ma^{-6},\qquad
 \E_Q n_{\mathrm{del}}^2\leq C_M\beta a^{-6}.
 \label{eq:patch-counts}
\end{equation}
For every candidate $t_0$, the total out count is bounded by the raw
count in $B(y,7a)$, which lies in a fixed annulus about the nucleus
with radius comparable to $a$.  The annular count estimate
\cref{eq:annular-count}, with \cref{eq:patch-offsets}, bounds the
expected square of this encompassing count by $C_Ma^{-6}$, uniformly
in $t_0$.  The deleted particles lie in the raw strip
$t_0-7b\leq|x_i-y|\leq t_0+b$.
Integrating its squared count over $t_0\in[5a,6a]$, each ordered pair
contributes for a set of lengths at most $8b$.  Averaging selects $t_0$
for the second bound in \cref{eq:patch-counts}, while preserving the
first bound.

\paragraph{The fresh conditional core.}
In the cut marginal of the tilted state, let $X$ record the out labels,
positions, and spins.  Slice only on these fresh cut data, not additionally
on the noisy arrays.  Denote the normalized core slice's density by
$\rho_X^c$ and its energy, including the price per core electron, by
$e_X^\lambda$.  Then
\[
 e_X^\lambda\geq P(\lambda),\qquad
 \Phi=V-\lambda-K*\rho_X^c.
\]
The localization products preserve antisymmetry within each group, and
the normalized core slices are in their sector form domains for almost
every datum \cite[\FLocalization]{Foundation}.

There is a nonnegative random variable $F=F(X)$ such that
\begin{equation}
 \Phi\leq F\text{ on }B(y,7a),\qquad
 \norm{F}_{L^2(Q)}\leq C_Ma^{-4}.
 \label{eq:patch-field-envelope}
\end{equation}
For this, condition the unshifted field obtained by subtracting only
electrons outside $B(y,20a)$; all those electrons belong to the core.
This field is harmonic on $B(y,10a)$ and bounds $\Phi$ above there.
At $z\in B(y,10a)$, the inclusion
$B(y,20a)\subset B(z,Aa_z)$ bounds it above by
$\E_Q[J_z\mid X]$.
The cut satisfies \cref{lem:local-field-input} uniformly: a bounded
number of comparable cells covers its out support, its gradient
coefficients obey $D_z^2\leq C\beta^{-2}$, and
$a_zm_z\geq ca^{-2}$, so that $D_z^2/(a_zm_z)\leq Ca^{1.6}\leq C$.
That lemma bounds the positive part in $L^2(Q)$ at each target $z$.
Submeans on balls of radius $a$, followed by Minkowski's integral
inequality, give the positive supremum bound in
\cref{eq:patch-field-envelope}.

\paragraph{The patch minimizer.}
On $\Omega=B(y,t_0-4b)$ minimize
\[
 \mathcal E_\Phi(f)=T\int f^{5/3}-\int\Phi f+D(f),\qquad f\geq0,
 \quad\supp f\subset\overline\Omega.
\]
The field $\Phi$ is bounded and harmonic on a neighborhood of
$\overline\Omega$: the core is excluded from $B(y,t_0)$, and the
nucleus is far outside this patch.  Coercivity in $L^{5/3}(\Omega)$,
weak lower semicontinuity, and strict convexity give a unique minimizer
$\rho=\rho_X$.  Its Euler equation is
\begin{equation}
 W=\Phi-K*\rho,\qquad \rho=kW_+^{3/2}\text{ on }\Omega.
 \label{eq:patch-euler}
\end{equation}
Nonnegative variations give the one-sided condition where $\rho=0$;
variations of either sign on positive density give the equality.
Comparison with zero and
$D(\rho)\geq(\int\rho)^2/(24a)$ show that
\begin{equation}
 \int\rho\leq CaF.
 \label{eq:patch-mass-bound}
\end{equation}
These minimizers can be chosen measurably in $X$.  One way to see this is
to note their continuous dependence in $L^{5/3}$ on a field in
$L^{5/2}(\Omega)$: coercivity bounds a converging sequence of minimizers;
weak limits minimize the limiting functional by lower semicontinuity and
comparison with its minimizer; uniqueness identifies the limit, and
convergence of the minimum values forces convergence of the
$L^{5/3}$ norms.  Uniform convexity gives strong convergence.  The
Coulomb form is a continuous positive quadratic form on
$L^{5/3}(\Omega)$, as used in this argument.

Let $\vartheta_b$ be the radial probability kernel from
\cite[\FNeumann]{Foundation}, obtained by averaging a cube kernel over
translations and rotations.  It is bounded by $b^{-3}$ and supported in
$B(0,\sqrt3b)$.  Define
\[
 \sigma=\sum_{i\ \mathrm{retained}}\vartheta_b(\,\cdot-x_i).
\]
Every such smear lies in $\Omega$ by the $7b$ retention margin.
We claim
\begin{equation}
 \E_Q\left[D(\sigma-\rho)+\int W_-\sigma\right]
       \leq C_Ma^{-7+.02},\qquad
 \int_{B(y,4a)}\rho_{\mathrm{tilt}}^{5/3}\leq C_Ma^{-7},
 \label{eq:patch-comparison}
\end{equation}
where $\rho_{\mathrm{tilt}}$ is the ordinary one-particle density of the
tilted state.  In particular, it is not the conditional patch density
$\rho$.

\paragraph{Lower and upper energy comparisons.}
Write $(o,c)$ for the fresh square cut and, for each out-label set
$I\subset\{1,\ldots,n\}$, set
\[
 \Psi_I=\prod_{i\in I}o(x_i)\prod_{i\notin I}c(x_i)\Psi_{A'},
 \qquad
 T_o=\frac12\sum_{I\subset\{1,\ldots,n\}}\sum_{i\in I}
                    \norm{\nabla_i\Psi_I}_2^2.
\]
Thus $T_o$ is the total out kinetic energy in the localized state.
Adding price to the exact localization identity of
\cite[\FLocalization]{Foundation} gives
\begin{equation}
 \langle\Psi_{A'},H_{Z,n}\Psi_{A'}\rangle+\lambda n
       +\operatorname{err}_{\mathrm{loc}}
 =\E_Qe_X^\lambda+T_o
       -\E_Q\sum_{i\ \mathrm{out}}\Phi(x_i)
       +\E_Q\sum_{i<l\ \mathrm{out}}K(x_i-x_l),
 \label{eq:priced-localization-identity}
\end{equation}
with $0\leq\operatorname{err}_{\mathrm{loc}}\leq C_Mb^{-2}a^{-3}$.
The price has been included in the core energy and in $\Phi$; there is
no localization error associated with it.  \cite[\FNeumann]{Foundation}
allows any measurably selected retained subset and gives
\begin{equation}
 \begin{split}
 T_o&\geq\E_Q\left[T\int\sigma^{5/3}
                              -Cb^{-2}(n_o^{4/3}+n_o)\right],\\
 \sum_{i<l\ \mathrm{out}}K(x_i-x_l)&\geq D(\sigma)-Cn_o/b.
 \end{split}
 \label{eq:patch-lower-comparison}
\end{equation}
The leading kinetic coefficient here is exactly $T$, for two spin states
and kinetic energy $-\Delta/2$.

For almost every fixed slice, the corresponding priced upper comparison is
\begin{equation}
 P(\lambda)\leq e_X^\lambda+T\int\rho^{5/3}
                   +Cb^{-2}\int\rho-\int\Phi\rho+D(\rho).
 \label{eq:patch-upper-comparison}
\end{equation}
The coherent-packet construction is standard; see
\cite[Lemma~8.2]{Solovej2003}. We spell out why the insertion proof of
\cite[\FInsertion]{Foundation} also proves this priced version.
Integrate the rank-one matrices of the orbitals
$g_b(x-z)e^{ip\cdot x}$, for both spins, over
$|p|\leq(3\pi^2\rho(z))^{1/3}$ with measure $dz\,dp/(2\pi)^3$.
The resulting one-body matrix $\gamma$ satisfies
\[
 0\leq\gamma\leq1,\qquad \Tr\gamma=\int\rho,\qquad
 \rho_\gamma=g_b^2*\rho,
\]
by Plancherel and momentum-ball integration.  Its kinetic trace is
\[
 \Tr(-\Delta/2)\gamma=T\int\rho^{5/3}+Cb^{-2}\int\rho.
\]
First truncate its spectral decomposition to finitely many eigenvectors.
Occupy these independently with probabilities given by the eigenvalues.
Every realization is a finite Slater determinant, with an integer number
of electrons, and its mean repulsion is at most the direct energy of the
truncated density by the direct-minus-exchange identity.  Its support is
in the packet region, which lies strictly inside $B(y,t_0)$ and is
separated from the core and nucleus.  It may therefore be wedged with
the core.  Every resulting state has shifted energy at least
$P(\lambda)$, and the mean price paid by the inserted electrons is
$\lambda\Tr\gamma$.

Positive spectral truncations have increasing densities and convergent
traces.  The potentials are bounded on the packet region for each fixed
slice, so all one-body terms converge; direct energies converge by
monotone convergence.  Radial smearing does not increase $D$, and
harmonic means replace $\int\Phi(g_b^2*\rho)$ by $\int\Phi\rho$,
including the constant price term.  Averaging the integer-sector trials
and passing to the limit proves \cref{eq:patch-upper-comparison}.
This reasoning requires neither an integer value of $\int\rho$ nor a
measurable choice of a trial determinant as the slice varies.

For retained particles, harmonic means likewise give
$\sum_{i\ \mathrm{retained}}\Phi(x_i)=\int\Phi\sigma$.
The deleted attractions cost at most $Fn_{\mathrm{del}}$.
Subtract \cref{eq:patch-upper-comparison} in
\cref{eq:priced-localization-identity}, use
\cref{eq:shifted-tilt-cost,eq:patch-lower-comparison}, and expand the
functional about its minimizer.  Convexity of the kinetic term leaves
$D(\sigma-\rho)$ and the linear term
\[
 \int\left(\frac{5T}{3}\rho^{2/3}-W\right)(\sigma-\rho)
                 =\int W_-\sigma.
\]
The equality uses \cref{eq:patch-euler}: where $\rho>0$ the coefficient
vanishes, and where $\rho=0$ it equals $W_-$.
By \cref{eq:patch-counts,eq:patch-field-envelope,eq:patch-mass-bound},
the resulting bound is
\begin{equation}
 \E_Q\left[D(\sigma-\rho)+\int W_-\sigma\right]
 \leq\delta'+C_M\left(
       \sqrt\beta\,a^{-7}+b^{-1}a^{-3}
                       +b^{-2}(a^{-4}+a^{-3})\right).
 \label{eq:patch-error-budget}
\end{equation}
All terms are integrable: for example, $\int\rho\leq CaF$ and
$\mathcal E_\Phi(\rho)\leq0$ bound both its kinetic and direct terms
by $CaF^2$ in expectation.  The exact cut identity also has integrable
absolute Coulomb terms by form-domain admissibility.

Since $b=a^{1.2}$, the errors in
\cref{eq:patch-error-budget}, divided by $a^{-7}$, are bounded by
\[
 C_M\left(a^{4.98}\log^5(e/a)+a^{.1}
                          +a^{2.8}+a^{.6}+a^{1.6}\right).
\]
This proves the first part of \cref{eq:patch-comparison}.
For the second, use $e_X^\lambda\geq P(\lambda)$ directly in
\cref{eq:priced-localization-identity}.  Dropping the nonnegative out
repulsion and using $\E_QFn_o\leq C_Ma^{-7}$ gives $T_o\leq C_Ma^{-7}$.
The coarse kinetic-density inequality of
\cite[\FLocalization]{Foundation}, on the full-out ball containing
$B(y,4a)$, proves the remaining assertion.

\subsection{Recovering the conditional density from the patch}

By \cref{eq:patch-euler}, the patch field satisfies
$\Delta W=4\pi kW_+^{3/2}$ on $B(y,4a)$.  We will use the bounds
\begin{equation}
 W\leq Ca^{-4}\text{ on }B(y,3a),\qquad
 \sup_{B(y,2t)}|W-W(y)|
                       \leq C\tau\bigl(a^{-4}+W(y)_-\bigr).
 \label{eq:patch-field-oscillation}
\end{equation}
For the first, compare on the ball of radius $R=4a$ with
\[
 U(x)=\frac{CR^4}{(R^2-|x-y|^2)^4}.
\]
For large universal $C$, $\Delta U\leq4\pi kU^{3/2}$; it diverges at
the boundary.  Testing with the positive part of $W-U$ proves the bound.
For the second, the Coulomb potential of
$\rho\ind_{B(y,2a)}$ is $O(a^{-4})$ with gradient $O(a^{-5})$ on that
ball, by the first bound and direct integration of the Coulomb kernel.
Adding this potential to $W$ leaves a harmonic function bounded above
by $Ca^{-4}$.  Apply the interior gradient estimate to its nonnegative
harmonic distance from this upper bound.  Its center value is at most
$Ca^{-4}+W(y)_-$, which proves the oscillation estimate.  These
comparisons are classical or weak positive-part tests: compact
$L^{5/3}$ densities have continuous potentials, so $W$ is continuous.

Define the local averaging operator
\[
 \mathcal R f(y)=\int\mathcal K_z(y)f(z)\,dz.
\]
On every compatible path under $Q$, including its fresh cut labels,
\begin{equation}
 |\mathcal R\sigma(y)-\mu_j(y)|
        \leq C_Mr^{-3}(\ell_j+b)t^{-4}=o(a^{-6}).
 \label{eq:noisy-array-sandwich}
\end{equation}
Here is the reason this comparison uses the original observation
posterior correctly.  Match any compatible raw configuration to the
$j$th observed multiset.  Each matched displacement is at most
$\sqrt3\ell_j$, so the raw kernel sum is within
$C_Mr^{-3}\ell_jt^{-4}$ of the sum evaluated at that multiset.
The same is true of its posterior average $\mu_j$ at these data.
Only $C_Mr^{-3}$ entries can contribute, by $\mathcal C_j$.
All relevant raw centers lie in the fully retained region of the fresh
cut, since $t/a\to0$, $b/a\to0$, and $\ell_j/a\to0$.
Replacing each such center by its fine smear changes the sum by at most
$C_Mr^{-3}bt^{-4}$.  This proves
\cref{eq:noisy-array-sandwich} without identifying the two posterior laws.
The ratios of these errors to $a^{-6}$ are bounded by
\[
 C_M\bigl(a^{.01-4w}+a^{.2-4w}\bigr)=o(1).
\]

Let
\[
 G=\{D(\sigma-\rho)\leq a^{-7+.01}\}.
\]
By \cref{eq:patch-comparison}, $Q(G^c)\leq C_Ma^{.01}$.
The function $z\mapsto\mathcal K_z(y)$ is a compactly supported
Lipschitz test with gradient norm at most $Ct^{-5/2}$.
Coulomb duality \cite[\FCoulomb]{Foundation} therefore gives, on $G$,
\begin{equation}
 |\mathcal R\rho(y)-\mathcal R\sigma(y)|
              \leq Ca^{-7/2+.005}t^{-5/2}=o(a^{-6}).
 \label{eq:local-coulomb-duality}
\end{equation}
Indeed the relative gain is at least $a^{.005-(5/2)w}$.
All the small errors here are uniform for fixed parameters along
\cref{eq:scaling-regime}, including scales with arbitrarily small $r/s$.

On $G$, the high-density antecedent in
\cref{eq:fixed-point-failures} forces
\begin{equation}
 W(y)\geq(h-\xi/16)d_y^{-4}.
 \label{eq:high-patch-field}
\end{equation}
Otherwise \cref{eq:patch-field-oscillation} would imply
$W(z)\leq(h-\xi/32)d_y^{-4}$ on the support of the averaging kernel.
This remains true for arbitrarily negative $W(y)$: its upper estimate
contains the nonpositive term $-(1-C\tau)W(y)_-$.
The Euler relation, \cref{eq:packet-approximate-mass}, and
\cref{eq:noisy-array-sandwich,eq:local-coulomb-duality} would then
contradict the high-density antecedent.  In the same way the low-density
antecedent forces, on $G$,
\begin{equation}
 W(y)\leq(h+\xi/16)d_y^{-4}.
 \label{eq:low-patch-field}
\end{equation}
The upper bound in \cref{eq:patch-field-oscillation} lets
\cref{eq:low-patch-field} pass to expectation with an $o(a^{-4})$ error.

The high case needs a bound on the negative field on $G^c$.
The high antecedent and \cref{eq:noisy-array-sandwich} hold on all paths
under $Q$, and give
\[
 \int_{B(y,2t)}\sigma\geq ca^{-6}t^3.
\]
Using \cref{eq:patch-field-oscillation} on this ball and absorbing
$C\tau W(y)_-$ yields
\begin{equation}
 W(y)_-\leq C(a^{-6}t^3)^{-1}\int W_-\sigma+C\tau a^{-4}.
 \label{eq:negative-patch-control}
\end{equation}
Consequently \cref{eq:patch-comparison} implies
\[
 \E_Q[\ind_{G^c}W(y)_-]=o(a^{-4}),
\]
since the ratio of its bound to $a^{-4}$ is at most
$C_Ma^{.02}\tau^{-3}+C\tau
 \leq C_Ma^{.02-3w}+Cs^w=o(1)$.
Thus \cref{eq:high-patch-field} also passes to expectation with
an $o(a^{-4})$ error.  This step explicitly controls the negative part;
no lower cap on the conditional field is assumed.

\subsection{Matching the potentials of the two conditional laws}

We adapt the averaged field transfer of
\cite[\FAveragedTransfer]{Foundation} to the present priced fields.
The preceding patch estimates supply its required small errors;
we give the two-law decomposition and estimates explicitly.
We prove
\begin{equation}
 \left|\E_Q\bigl(H'_j-W\bigr)(y)\right|=o(a^{-4}).
 \label{eq:potential-matching}
\end{equation}
Write $P^{\mathrm{raw}}$ and $P^{\mathrm{ms}}$ for the raw and
master-smeared electron potentials at $y$, and
$P^{\mathrm{ret}},P^{\mathrm{del}}$ for the raw potentials of retained
and deleted out particles.  Because the event defining $Q$ is in
$\mathcal F_j$, the original observation tower gives
$\E_Q(K*\mu_j)(y)=\E_QP^{\mathrm{ms}}$.
The core in $W$ instead uses the fresh-cut tower in the tilted raw law.
Combining these two identities, and cancelling both nucleus and price,
expresses $\E_Q(H'_j-W)(y)$ as the expectation of
\begin{equation}
 (P^{\mathrm{raw}}-P^{\mathrm{ms}})-P^{\mathrm{del}}
    -(P^{\mathrm{ret}}-K*\sigma)(y)+K*(\rho-\sigma)(y).
 \label{eq:two-posterior-potential-identity}
\end{equation}

By H\"older's inequality and the ordinary density bound in
\cref{eq:patch-comparison}, for $0<u\leq4a$ the expected raw potential
from $B(y,u)$ is at most
\begin{equation}
 \int_{B(y,u)}\frac{\rho_{\mathrm{tilt}}(z)}{|z-y|}\,dz
                  \leq C_Ma^{-4}(u/a)^{1/5}.
 \label{eq:raw-short-range-potential}
\end{equation}
Radial smearing never increases the potential of a point charge and
leaves it unchanged outside the smear radius.  The master-smearing
difference in \cref{eq:two-posterior-potential-identity} can therefore
come only from centers within $2t$ of $y$; the fine-smearing difference
can come only from centers within $\sqrt3b$.
\Cref{eq:raw-short-range-potential} bounds both by $o(a^{-4})$.
Deleted centers have distance comparable to $a$, so
\cref{eq:patch-counts} gives
$\E_QP^{\mathrm{del}}\leq C_Ma^{-4}\sqrt\beta$.

For the remaining signed potential, truncate the kernel at height $t^{-1}$.
The test $z\mapsto\min(|z-y|^{-1},t^{-1})$ has homogeneous Sobolev gradient
norm $O(t^{-1/2})$.  Coulomb duality, or its compactly supported
approximation, and \cref{eq:patch-comparison} give expected absolute
truncated error at most
\[
 C_Mt^{-1/2}a^{-7/2+.01}=o(a^{-4});
\]
the relative gain is $a^{.01}\tau^{-1/2}\leq C_Ma^{.01-w/2}$.
The truncation loss for $\rho$ is at most $Ca^{-6}t^2$ by
\cref{eq:patch-field-oscillation}.  For $\sigma$, radial smearing and
\cref{eq:raw-short-range-potential} bound it by
\[
 C_Ma^{-4}\bigl((t+\sqrt3b)/a\bigr)^{1/5}.
\]
These also vanish relative to $a^{-4}$ by
\cref{eq:packet-scale-ratios}.  This proves
\cref{eq:potential-matching}.

In the high case, \cref{eq:high-patch-field,eq:negative-patch-control}
give $\E_QW(y)\geq(h-\xi/16)d_y^{-4}-o(a^{-4})$, whereas the failed
comparison in \cref{eq:fixed-point-failures} gives
$\E_QH'_j(y)\leq(h-\xi/4)d_y^{-4}$.
\Cref{eq:potential-matching} contradicts these inequalities
eventually.  The low case follows in the same way from
\cref{eq:low-patch-field} and the upper bound on $W$.
We conclude that each fixed-point failure in
\cref{eq:fixed-point-failures}, intersected with $\mathcal C_j$, has
probability less than $r^{42}$.

\subsection{Spatial and height nets}

Take a spatial net of mesh at most $r^2$ in the original band, with
$O_M(r^{-3})$ points, and a fixed finite height net in
$[h_l/2,2h_h]$ of mesh at most $\xi/64$.
Exclude the fixed-point failures at all net pairs, the complement of
$\mathcal C_j$, and the enlarged-cap failure.  Their total probability
is $O_M(r^{39})+O_M(r^{40})$, in particular at most $Cr^{25}$.
All these events are $\mathcal F_j$-measurable.

The normalized density $d_y^6\mu_j(y)$ changes by at most
$C_Mr^{1-4w}=o(1)$ between a point and its nearest net point, by
\cref{eq:conditional-density-bounds}.  For the field, work on balls of
radius a small fixed multiple of $r$ within the enlarged cap band.
Adding the Coulomb potential of the local source $\mu_j$ makes $H'_j$
harmonic there and bounds that harmonic function above by
$C_Mr^{-4-3w}$, using
\cref{eq:conditional-density-bounds,eq:enlarged-field-cap}.
The positive-harmonic gradient argument used in
\cref{eq:patch-field-oscillation} shows that
$v(y)=d_y^4H'_j(y)$ varies over a displacement $O(r^2)$ by at most
\begin{equation}
 C_Mr^{1-3w}+C_Mr\,v(y)_-.
 \label{eq:net-field-variation}
\end{equation}

Suppose a high-density failure with tolerance $\xi$ occurs at $(y,h)$.
If $v(y)\geq0$, \cref{eq:net-field-variation} puts the nearby net value
below $h-3\xi/4$ for small $s$.  If $v(y)<0$, the same upper estimate
is bounded by
$(1-C_Mr)v(y)+C_Mr^{1-3w}$, and gives the identical conclusion since
$h\geq h_l>16\xi$.  Choose a height-net value between
$h-\xi/4$ and $h-\xi/8$.  The small normalized-density variation
preserves the strict high antecedent at this lower height, while the
field value violates its comparison with tolerance $\xi/4$.
This is an excluded net failure.

A low-density failure has $v(y)>h+\xi>0$.  Its nearest net value is above
$h+3\xi/4$ by \cref{eq:net-field-variation}.  A height-net value between
$h+\xi/8$ and $h+\xi/4$ preserves the low antecedent and again yields
an excluded failure with tolerance $\xi/4$.
Thus both comparisons hold everywhere on the band for every
$h_l\leq h\leq h_h$.  Together with
\cref{eq:enlarged-field-cap}, this proves
\cref{prop:inverse-comparison}.
